\section{Parametric description of planes}

The parametric description of a plane is based on movement. A point moves in two independent directions. This is the main difference between the parametrization of a line and the parametrization of a plane.

A line needs one parameter. A plane needs two parameters.

\subsection*{Vector parametric form}

Let \(P_0\) be a point in space and let \(u\) and \(v\) be two non-parallel vectors. The plane through \(P_0\) in the directions \(u\) and \(v\) is
\[
P=P_0+s u+t v,
\qquad s,t\in\mathbb{R}.
\]
The parameters \(s\) and \(t\) move independently.

\begin{examplebox}
Consider the plane
\[
P=(2,-1,3)+s(1,0,2)+t(0,3,-1).
\]
For \(s=0\), \(t=0\), we get
\[
P=(2,-1,3).
\]
For \(s=1\), \(t=0\), we move by \((1,0,2)\):
\[
P=(3,-1,5).
\]
For \(s=0\), \(t=1\), we move by \((0,3,-1)\):
\[
P=(2,2,2).
\]
All these points lie on the same plane.
\end{examplebox}

The two direction vectors should not be parallel. If they are parallel, the two-parameter formula does not sweep out a plane; it only moves along one line.

\subsection*{Coordinate parametric form}

If
\[
P_0=(x_0,y_0,z_0),
\qquad
u=(a_1,a_2,a_3),
\qquad
v=(b_1,b_2,b_3),
\]
then
\[
(x,y,z)=(x_0,y_0,z_0)+s(a_1,a_2,a_3)+t(b_1,b_2,b_3).
\]
In coordinate form,
\[
x=x_0+sa_1+tb_1,
\]
\[
y=y_0+sa_2+tb_2,
\]
\[
z=z_0+sa_3+tb_3.
\]
The same values of \(s\) and \(t\) must be used in all three coordinates.

\subsection*{Checking whether a point lies on a parametric plane}

To check whether a point lies on a parametric plane, we ask whether there exist values of the two parameters that produce the point.

\begin{examplebox}
Consider
\[
(x,y,z)=(1,0,2)+s(1,1,0)+t(0,2,1).
\]
Does the point \((3,4,3)\) lie on the plane?

The coordinate equations are
\[
x=1+s,
\qquad
 y=s+2t,
\qquad
 z=2+t.
\]
From \(x=3\), we get
\[
s=2.
\]
From \(z=3\), we get
\[
t=1.
\]
Then the middle coordinate gives
\[
y=s+2t=2+2=4.
\]
Therefore \((3,4,3)\) lies on the plane.
\end{examplebox}

A point belongs to the plane only if all coordinate equations are satisfied by the same pair \((s,t)\).

\subsection*{From parametric form to normal form}

If a plane is given parametrically by
\[
P=P_0+s u+t v,
\]
then a normal vector can be found by
\[
n=u\times v.
\]
Then the normal equation is
\[
n\cdot(P-P_0)=0.
\]

This conversion is useful because distances and angles are usually easier to compute from normal vectors.

\begin{summarybox}
When reading a parametric plane, ask:
\begin{itemize}
\item What point anchors the plane?
\item What are the two direction vectors?
\item Are the direction vectors non-parallel?
\item What do the two parameters control?
\item Can I find a normal vector using a cross product?
\item Does a given point arise from one common pair of parameter values?
\end{itemize}
\end{summarybox}

The parametric form describes a plane by motion. It is the language of two independent directions.